\documentclass[11pt,a4paper]{article}
\usepackage[margin=2.2cm]{geometry}
\usepackage{amsmath,amssymb}
\usepackage{booktabs}
\usepackage{textcomp}
\usepackage{xcolor}
\usepackage[hidelinks]{hyperref}
\newcommand{\E}{\mathbb{E}}
\newcommand{\Prob}{\mathbb{P}}
\newcommand{\meur}{M\texteuro}
\newcommand{\code}[1]{\texttt{#1}}

\title{Lending Game -- Stage 1 Strategy Note\\[2pt]
\large Team ``Lender 2'' -- Introduction to Finance for Data Scientists, HEC Paris}
\author{Internal note for the team's review before the October 17 submission}
\date{October 13, 2026}

\begin{document}
\maketitle

\section*{Summary}
\begin{itemize}
  \item We are lender 2. For the 100,000 new applicants we observe sex, employment, marital status, income and our
        private signal $s_2$; the competitors observe $s_1$ or $s_3$ instead. The past-loans file contains all
        three signals and the outcome, which lets us \emph{measure} how much the competitors know that we do not,
        and \emph{simulate} the whole market before submitting anything.
  \item Default is driven mostly by income and employment. Each private signal adds a little (AUC $0.717 \to 0.735$),
        the three together add more ($0.760$). The resulting hidden information has a standard deviation of
        4.5 points of default probability -- larger than any sensible margin -- and is concentrated at low income
        (7.1 points at zero income, 0.3 points above 40k).
  \item Applicants take the cheapest offer, so we win a loan exactly when the competitors' signals were bad: the
        winner's curse. Naive pricing (break-even $+3\%$) earns 1.7~\meur{} on average over twelve competitor
        scenarios and loses up to 2.2~\meur{}.
  \item Our rule: break-even plus a 3--4\% margin when our estimated PD is below an income-dependent threshold
        (8--12\% below 20k of income, 25\% above); otherwise an 82.8\% ``lender of last resort'' offer, which is
        only accepted when nobody else lends cheaper. Average 6.4~\meur{} over the twelve scenarios, worst
        scenario $+1.2$~\meur{}, market share about 22\%, first in every scenario.
  \item The level of that last-resort rate is the decisive choice of stage 1. If another team also uses a safety
        net with low thresholds, the 10--25\% PD band (34\% of applicants) becomes a price war worth 65 to
        190~\meur{} to whoever is lowest within 2\%, against 5~\meur{} for the whole normal game. We chose 0.828:
        it wins that war against any competitor net at 0.85 or above and costs only 0.35~\meur{} if nobody nets.
  \item The submission contains an offer for every applicant: 54,418 normal offers between 3.0\% and 17.6\%,
        and 45,582 offers at 82.8\%.
\end{itemize}

\section{The game and what we optimize}
Each of 100,000 applicants asks for a 10,000~\texteuro{} loan from three competing lenders. Lender $k$ offers a rate
$r(k,i)$ or nothing. One third of the applicants prefers each lender and is willing to pay 2\% more to borrow from
it: a type-$k$ applicant takes the lowest of the three rates after subtracting $0.02$ from lender $k$'s rate.
If the loan is repaid we earn $10{,}000\,r$; if the borrower defaults we lose $10{,}000$ (zero recovery).
Our score is total profit, judged in level and in rank against the two other teams (10\% of the grade for
stage~1, 20\% for stage~2; the report is 70\%).

With tens of thousands of loans, sampling noise is small (the standard deviation of total profit is about
0.5~\meur{} for 30,000 loans). The real uncertainty is strategic: what the other teams do. The stage-1 goal is
therefore a strategy that is profitable against any plausible competitor, never loses money, and lets us learn
the competitors' behaviour for stage~2, where the aggressive moves belong.

\section{What the data says}
Past loans: 100,000 rows, 9.95\% defaults, no missing values; the new applications match the past population on
every moment we checked (shares of employment status, income quantiles, signal distribution). Table~\ref{tab:data}
shows the main determinants of default.

\begin{table}[h]\centering\small
\begin{tabular}{lrr@{\hskip 1.2cm}lrr}
\toprule
Income & Share & Default & Other variable & Share & Default\\
\midrule
0           & 17.0\% & 20.2\% & unemployed & 14.3\% & 15.6\%\\
1--5k       &  6.1\% & 15.1\% & student    & 16.2\% & 12.9\%\\
5--10k      & 29.9\% & 12.5\% & retired    & 11.6\% &  9.0\%\\
10--20k     & 25.9\% &  5.8\% & employed   & 57.9\% &  7.9\%\\
20--40k     & 12.8\% &  2.2\% & $s_2 = 0$ (``no signal'') & 20.9\% & 9.7\%\\
40--80k     &  5.1\% &  0.8\% & $s_2 \in (0, 0.2]$   & 12.9\% & 14.9\%\\
80k+        &  3.1\% &  0.1\% & $s_2 \in (0.8, 1]$   & 15.4\% &  7.9\%\\
\bottomrule
\end{tabular}
\caption{Default rates in the past loans. Women default slightly more than men (10.7\% vs 9.4\%) and singles more
than married applicants (10.8\% vs 9.4\%).}
\label{tab:data}
\end{table}

Three facts matter for the strategy.
\begin{itemize}
  \item \textbf{Income dominates.} Default goes from 20\% at zero income to 0.1\% above 80k. The signals are
        weaker and non-linear: 15\% default in the lowest bucket, flat around 8\% above $0.6$.
  \item \textbf{A signal of exactly zero means ``no signal'', not ``worst score''.} 21\% of each signal is exactly
        0, and those applicants default at the population average. A model that treats 0 as a very low score
        overprices a fifth of the applicants. Tree models isolate the spike on their own; we also add an explicit
        indicator.
  \item \textbf{The three signals are different information.} Their pairwise correlation is only $0.18$ and the
        zeros are independent across signals. Hence the competitors genuinely know things about each applicant
        that we do not -- the source of adverse selection in this game.
\end{itemize}

\begin{table}[h]\centering\small
\begin{tabular}{lcc@{\hskip 1.2cm}lc}
\toprule
Model (out of sample, 5-fold) & AUC & log-loss & Income segment & $\sigma(U)$ \\
\midrule
$X$ only (sex, employment, married, income) & 0.717 & 0.2951 & 0       & 0.071\\
$X + s_2$ (us)                              & 0.735 & 0.2908 & 1--5k   & 0.050\\
$X + s_1$ (lender 1)                        & 0.735 & 0.2908 & 5--10k  & 0.052\\
$X + s_3$ (lender 3)                        & 0.736 & 0.2904 & 10--20k & 0.025\\
$X + s_1 + s_2 + s_3$ (full information)    & 0.760 & 0.2828 & 20--40k & 0.008\\
                                            &       &        & 40k+    & 0.003\\
\bottomrule
\end{tabular}
\caption{Left: predictive power of the information sets. Right: standard deviation of the hidden information
$U = p_{\text{full}} - \hat p_2$, i.e.\ how much the full-information default probability deviates from our own
estimate. Overall $\sigma(U) = 0.045$.}
\label{tab:models}
\end{table}

Table~\ref{tab:models} quantifies the asymmetry. Our model and the competitors' are equally good (same data, same
method, symmetric signals), but a lender who saw all three signals would correct our estimate by $\pm 4.5$ points
of PD on average. At high income the competitors' signals are nearly worthless ($\sigma(U) = 0.3$ points); at zero
income they move the PD by 7 points. This map tells us where we can compete and where we cannot.

\section{The economics behind the pricing rule}

\subsection{Break-even rate}
Interest is received only if the borrower repays, so the expected profit of a loan at rate $r$ to a borrower with
default probability $p$ is $10{,}000\,[(1-p)\,r - p]$ and the break-even rate is
\begin{equation}
  r^*(p) = \frac{p}{1-p}, \qquad\text{not } r = p.
  \label{eq:breakeven}
\end{equation}
The lecture formula $\text{NPV} = R - \Prob[\text{default}]$ assumes the interest is paid at issuance. A team
applying it literally and pricing $r = p + m$ offers 23\% to a 20\%-PD applicant whose break-even is 25\%; it wins
those loans and loses money on them. We keep that behaviour as one of the competitor archetypes.

\subsection{The winner's curse through the competitors' signals}
Let $\hat p_2 = \E[D \mid X, s_2]$ be our estimate and $p = \E[D \mid X, s_1, s_2, s_3] = \hat p_2 + U$ the
full-information probability, with $\E[U \mid X, s_2] = 0$. Our estimate is unbiased (Table~\ref{tab:calib}), but
the applicants who \emph{accept} our offer are not a random sample: an applicant takes our offer when the
competitors' effective rates are above ours, and competitors price higher when their signals are bad, i.e.\ when
$U$ is high. Writing $Q_i(r)$ for the probability that applicant $i$ accepts our offer at rate $r$ and
$\pi_i(r) = \E[D_i \mid \text{accepts at } r]$ for the default rate of the accepting borrowers, the expected
profit of the offer is
\begin{equation}
  \Pi_i(r) = 10{,}000 \cdot Q_i(r)\,\big[(1-\pi_i(r))\,r - \pi_i(r)\big], \qquad \pi_i(r) > \hat p_{2,i}.
  \label{eq:profit}
\end{equation}
This is the lecture's $M - \E[U \mid U > M - \bar R] < M$ with the acceptance event generated by the competitors'
information rather than by the borrower's. Both $Q_i$ and $\pi_i$ depend on what the competitors do, which is why
we simulate them instead of adding a fixed margin. The curse is measurable on the past loans: with naive
break-even $+3\%$ pricing against competitors doing the same, the loans we win default at 10.0\% while we
predicted 8.1\%; pricing \emph{higher} makes it worse (12.1\% realized vs 8.7\% predicted at $+5\%$), because a
high price only wins the applicants everyone else found risky.

\subsection{The 2\% preference: fans versus everyone}
For an applicant whose best competitor rate is $c$, our demand is a step function: we win our own fans (one third
of applicants) if $r < c + 0.02$, and all three types if $r < c - 0.02$. Undercutting by 4 points therefore
triples volume. Ignoring selection, the all-market price beats the fans-only price when
$3[(1-p)(r-0.04) - p] > (1-p)\,r - p$, i.e.\ when
\begin{equation}
  r > \frac{p}{1-p} + 0.06 .
\end{equation}
Grabbing market share pays only when the competitors sit more than 4 points above break-even, which teams
equipped with a decent PD model will not do. The natural target in the safe segments is therefore our fans, at
break-even plus a few points. Simulation confirms it: pricing at break-even $+1\%$ to gain share gives a 55\%
market share and a loss of 2~\meur{}, and drags the competitors into losses as well.

\subsection{The lender of last resort}
Applicants have no outside option: whoever receives at least one offer borrows. An offer at a very high rate
$r_{\text{net}}$ is accepted only when every competitor's effective rate is above ours or absent, so it costs
nothing when the competitors lend, and when both reject the applicant we lend at $r_{\text{net}}$ to a borrower with
default probability $d = \E[D \mid \text{both reject}]$. The loan is profitable if
\begin{equation}
  (1-d)\,r_{\text{net}} - d > 0 \iff r_{\text{net}} > \frac{d}{1-d}.
  \label{eq:net}
\end{equation}
Table~\ref{tab:pool} evaluates the abandoned pool on the past loans for competitors who reject above a PD
threshold: as long as they reject at 25\% or below, any rate above 74\% is profitable; if they reject only above
30\% the pool is unprofitable at any rate but represents 0.2\% of applicants (a loss of about 0.05~\meur{} at
worst). This ``safety net'' replaces ``no offer'' everywhere in our submission. Its \emph{level} matters only
against other lenders' nets, and Section~\ref{sec:net} shows that this contest dwarfs everything else in the game.

\begin{table}[h]\centering\footnotesize
\begin{tabular}{crrrrrrr}
\toprule
Competitors reject if PD $>$ & Pool & Default $d$ & Break-even $d/(1-d)$ & \multicolumn{4}{c}{Profit per loan (\texteuro) at a net of}\\
 & & & & 0.989 & 0.90 & 0.80 & 0.70\\
\midrule
15\% & 14.4\% & 24.3\% & 32\% & 5,061 & 4,387 & 3,630 & 2,873\\
20\% &  5.3\% & 29.9\% & 43\% & 3,934 & 3,311 & 2,610 & 1,910\\
25\% &  0.8\% & 42.6\% & 74\% & 1,420 &   909 &   335 &  $-240$\\
30\% &  0.2\% & 51.2\% & none & $-289$ & $-724$ & $-1,212$ & $-1,700$\\
\bottomrule
\end{tabular}
\caption{The applicants both competitors would reject (their own PD estimate above the threshold), on the past loans.}
\label{tab:pool}
\end{table}

\section{Methodology}

\subsection{Default probability models}
Gradient-boosted trees (\code{HistGradientBoostingClassifier}) on employment dummies, marital status, income,
$\log(1+\text{income})$, the signal and a ``signal $=0$'' indicator. Hyper-parameters were chosen by 3-fold
cross-validated log-loss (the right metric, since we need calibrated probabilities, not just ranking):
learning rate 0.05, 7 leaves, 200 iterations, at least 200 observations per leaf, $L_2 = 1$. Four models are
fitted with the same recipe: ours ($X + s_2$), each competitor's ($X + s_1$, $X + s_3$) and the full-information
model ($X + s_1 + s_2 + s_3$). All predictions used in the simulations are 5-fold out-of-sample, so no lender is
clairvoyant in the simulator; the model used to price the new applications is refitted on all 100,000 past loans.

\begin{table}[h]\centering\small
\begin{tabular}{lcccccccccc}
\toprule
Decile of $\hat p_2$ & 1 & 2 & 3 & 4 & 5 & 6 & 7 & 8 & 9 & 10\\
\midrule
Predicted & 0.007 & 0.022 & 0.035 & 0.054 & 0.077 & 0.101 & 0.123 & 0.145 & 0.179 & 0.253\\
Realized  & 0.006 & 0.021 & 0.033 & 0.053 & 0.077 & 0.102 & 0.122 & 0.143 & 0.182 & 0.256\\
\bottomrule
\end{tabular}
\caption{Calibration of our model (out of sample).}
\label{tab:calib}
\end{table}

Sex is excluded from the final model: with or without it the out-of-sample AUC is 0.7349 and the log-loss
0.2908. It adds nothing, and pricing on sex would be illegal for a real lender, so we drop it at zero cost.

\subsection{Market simulator}
\code{lending\_game/simulator.py} implements the rules literally: applicant types are drawn uniformly at random
(the rules state one third each), the 2\% preference is subtracted from the preferred lender's rate, the lowest
effective rate wins, ties are broken at random, and the profit is computed from the realized defaults of the past
loans. We also compute an ``expected'' profit using the full-information PD instead of the realized outcome: it has
the same mean (the full-information model is calibrated) but much less noise, so it is the objective of the
optimizer, while the realized profit is reported as a check. Every result is averaged over three independent
type draws.

\subsection{Competitor scenarios}
We expect the other teams to work with the same data and similar boosted-tree models, so no edge can come from the
model itself; the uncertainty is in their pricing rule. Each archetype in Table~\ref{tab:archetypes} maps a
competitor's own PD estimate to a rate, and twelve scenarios combine them: eight symmetric
(\code{base/base}, \code{aggressive/aggressive}, \code{cautious/cautious}, \code{no\_reject/no\_reject},
\code{slide/slide}, \code{slide\_norej/slide\_norej}, \code{smart/smart}, \code{base\_net/base\_net}) and four mixed
(\code{base/aggressive}, \code{cautious/slide}, \code{smart/base}, \code{loose/tight}). The \code{clone}
archetype -- a team that reasons exactly like us -- is studied separately in Section~\ref{sec:net}.

\begin{table}[h]\centering\small
\begin{tabular}{llll}
\toprule
Archetype & Rate & Rejects if PD $>$ & Instead of ``no offer''\\
\midrule
base        & $p/(1-p) + 3\%$ & 25\% & --\\
aggressive  & $p/(1-p) + 1\%$ & never & --\\
cautious    & $p/(1-p) + 5\%$ & 20\% & --\\
no\_reject  & $p/(1-p) + 3\%$ & never & --\\
loose / tight & $p/(1-p) + 2\%$ / $+4\%$ & 30\% / 20\% & --\\
slide / slide\_norej & $p + 3\%$ (lecture formula) & 25\% / never & --\\
smart       & $p/(1-p)$ + 2--6\% rising as income falls & 25\% & 0.99 safety net\\
base\_net   & $p/(1-p) + 3\%$ & 25\% & 0.99 safety net\\
clone       & our own rule (Table~\ref{tab:rule}) & 8--25\% as ours & net of its choosing\\
\bottomrule
\end{tabular}
\caption{Competitor archetypes, each applied to the competitor's own out-of-sample PD estimate.}
\label{tab:archetypes}
\end{table}

\subsection{Our pricing rule and its optimization}
For applicant $i$ with our PD estimate $\hat p_i$ and income segment $s_i$,
\begin{equation}
  r_i = \begin{cases}
     \dfrac{\hat p_i}{1-\hat p_i} + m(s_i) & \text{if } \hat p_i \le \tau(s_i),\\[8pt]
     r_{\text{net}} & \text{otherwise.}
  \end{cases}
  \label{eq:rule}
\end{equation}
The parameters are a margin $m$ and a threshold $\tau$ per income segment (six segments: 0, 1--5k, 5--10k,
10--20k, 20--40k, 40k+) and the net rate. They are chosen by coordinate descent on grids (margins 0.5--10\%,
thresholds 4--100\%, net in $\{\text{none}, 0.5, \dots, 0.989\}$) to maximize the expected profit averaged over the
twelve scenarios and three type draws, evaluated on the past loans. A richer rule with margins by segment
$\times$ risk tercile improved the objective by only 0.2\% and was discarded. The net rate was then fixed by the
dedicated analysis of Section~\ref{sec:net}.

\section{Results}

\subsection{The chosen rule}
\begin{table}[h]\centering\small
\begin{tabular}{lrrrrrr}
\toprule
Segment & New applicants & Mean $\hat p$ & Margin $m$ & Threshold $\tau$ & Offered the net & Median normal rate\\
\midrule
0       & 17,159 & 0.201 & 4.0\% & 12\% & 92\% & 17.4\%\\
1--5k   &  6,142 & 0.150 & 3.5\% & 12\% & 74\% & 15.7\%\\
5--10k  & 29,870 & 0.126 & 3.5\% & 10\% & 69\% & 12.3\%\\
10--20k & 25,766 & 0.058 & 3.0\% &  8\% & 18\% &  8.0\%\\
20--40k & 12,829 & 0.023 & 3.0\% & 25\% &  0\% &  5.3\%\\
40k+    &  8,234 & 0.006 & 3.0\% & 25\% &  0\% &  3.4\%\\
\bottomrule
\end{tabular}
\caption{Optimized parameters and their effect on the new applications; $r_{\text{net}} = 0.828$.}
\label{tab:rule}
\end{table}

The optimizer's answer is simple to state: \emph{compete normally where the competitors cannot know more than
us; where they can, either take the clearly safe applicants or be the lender of last resort -- never in
between.} Above 20k of income the thresholds are high and we lend to almost everyone at break-even $+3\%$. Below
20k, any applicant whose PD exceeds 8--12\% gets the net. The margins are not lower at high income because, with
the 2\% preference, cutting the margin from 3\% to 1\% does not win any non-fan and loses two thirds of the
profit on the fans (Section~3.3).

\subsection{Comparison with naive strategies}
\begin{table}[h]\centering\small
\begin{tabular}{lrrr}
\toprule
Strategy (realized profit on the past loans, 12 scenarios) & Mean & Worst scenario & Mean share\\
\midrule
Break-even $+3\%$, lend to everyone                 & 1.7~\meur{} & $-2.2$~\meur{} & 32\%\\
Break-even $+3\%$, no offer if PD $> 25\%$           & 2.3~\meur{} & $-1.9$~\meur{} & 32\%\\
Break-even $+3\%$, net at 0.989 instead of no offer & 2.9~\meur{} & $-1.9$~\meur{} & 32\%\\
\textbf{Optimized rule (net 0.828)}                  & \textbf{6.4~\meur{}} & $\mathbf{+1.2}$~\meur{} & 22\%\\
\bottomrule
\end{tabular}
\caption{Each strategy against the same twelve competitor scenarios.}
\label{tab:strategies}
\end{table}

\subsection{Scenario by scenario}
\begin{table}[h]\centering\footnotesize
\begin{tabular}{lrrrrrr}
\toprule
Scenario (lender 1 / lender 3) & Our share & Our profit & $\hat p_2$ on won & Realized default & Lender 1 & Lender 3\\
\midrule
base / base                 & 21.7\% &  5.1 & 0.054 & 0.068 &  3.2 &  3.2\\
aggressive / aggressive     &  9.1\% &  1.2 & 0.049 & 0.064 & $-4.0$ & $-4.0$\\
cautious / cautious         & 42.7\% & 24.8 & 0.066 & 0.076 &  6.5 &  6.2\\
no\_reject / no\_reject     & 20.8\% &  4.7 & 0.046 & 0.053 &  2.3 &  2.5\\
slide / slide               & 19.4\% &  4.4 & 0.051 & 0.067 & $-1.1$ & $-0.9$\\
slide\_norej / slide\_norej & 18.6\% &  4.0 & 0.043 & 0.050 & $-2.2$ & $-1.9$\\
smart / smart               & 24.7\% &  6.4 & 0.057 & 0.067 &  6.2 &  6.3\\
base\_net / base\_net       & 21.7\% &  5.1 & 0.054 & 0.068 &  3.2 &  3.2\\
base / aggressive           & 12.6\% &  2.2 & 0.050 & 0.061 & $-1.1$ & $-2.1$\\
cautious / slide            & 27.2\% &  8.7 & 0.056 & 0.071 &  0.3 &  2.7\\
smart / base                & 23.0\% &  5.6 & 0.055 & 0.068 &  3.3 &  5.5\\
loose / tight               & 20.0\% &  5.0 & 0.054 & 0.067 &  2.6 &  1.9\\
\bottomrule
\end{tabular}
\caption{Optimized rule: realized profits in \meur{}, averaged over three type draws.}
\label{tab:scenarios}
\end{table}

We are profitable in every scenario and first in every one of them; the closest contests are \code{smart/smart}
(6.4~\meur{} against 6.2--6.3~\meur{}, two competitors who reason like us above a 25\% threshold and keep their
own 10--25\% PD business at fat margins) and \code{smart/base} (5.6 against 5.5). When the competitors are
aggressive they lose 4~\meur{} each while we keep $+1.2$~\meur{}. The curse is tamed but visible: the loans we win
default about 1.4 points more than we predict (6.8\% vs 5.4\% in the base scenario), against 2 points with naive
pricing. The 25~\meur{} of the cautious scenario is the net at work on a 5\% abandoned pool; we do not count on
it.

\subsection{Where the profit comes from}
\begin{table}[h]\centering\small
\begin{tabular}{lrrrrrr}
\toprule
Segment & Applicants & Share won & Mean rate of won loans & $\hat p_2$ on won & Realized default & Profit\\
\midrule
0       & 17,035 &  7.9\% & 79.6\% & 0.180 & 0.298 & 0.27~\meur{}\\
1--5k   &  6,125 & 12.5\% & 67.8\% & 0.106 & 0.112 & 0.27~\meur{}\\
5--10k  & 29,904 & 14.4\% & 60.8\% & 0.080 & 0.100 & 0.68~\meur{}\\
10--20k & 25,923 & 31.9\% & 21.7\% & 0.046 & 0.053 & 1.73~\meur{}\\
20--40k & 12,780 & 33.9\% &  5.4\% & 0.022 & 0.023 & 1.24~\meur{}\\
40k+    &  8,233 & 33.2\% &  3.6\% & 0.006 & 0.007 & 0.80~\meur{}\\
\bottomrule
\end{tabular}
\caption{Optimized rule in the base scenario (one type draw). Above 10k of income we win one third of the
applicants -- our fans -- with no curse to speak of; below 10k most won loans are net loans to abandoned
applicants, which default at 30\% but are priced at 82.8\%.}
\label{tab:segments}
\end{table}

\subsection{Robustness of the thresholds}
Raising the thresholds to ``look more moderate'' is expensive (Table~\ref{tab:tau}): at $\tau = 25\%$ everywhere
the mean falls from 6.6 to 3.1~\meur{} and the worst scenario becomes a loss of 1.1~\meur{}. The low thresholds
are not an artefact of the scenario where both competitors reject a lot: re-optimizing without that scenario, only
against competitors who never reject anyone, or only against rejecting competitors, returns \emph{identical}
thresholds. What drives them is the curse itself: at low income, the 10--25\% PD band is where the competitors'
hidden information is largest, and competing for it at any normal rate loses money. Removing the net altogether
(plain ``no offer'') keeps the worst case but costs 2.0~\meur{} on average.

\begin{table}[h]\centering\small
\begin{tabular}{lrrr}
\toprule
Variant (same margins, net 0.828) & Mean expected profit & Worst scenario & Base scenario\\
\midrule
Thresholds as chosen       & 6.6~\meur{} & $+1.2$~\meur{} & 5.3~\meur{}\\
$\tau = 15\%$ everywhere   & 5.9~\meur{} & $+0.2$~\meur{} & 4.7~\meur{}\\
$\tau = 20\%$ everywhere   & 4.7~\meur{} & $-0.5$~\meur{} & 3.6~\meur{}\\
$\tau = 25\%$ everywhere   & 3.1~\meur{} & $-1.1$~\meur{} & 2.6~\meur{}\\
$\tau = 30\%$ everywhere   & 2.4~\meur{} & $-1.3$~\meur{} & 2.1~\meur{}\\
Chosen thresholds, no offer instead of the net & 4.6~\meur{} & $+1.2$~\meur{} & 4.6~\meur{}\\
\bottomrule
\end{tabular}
\caption{Sensitivity to the rejection threshold.}
\label{tab:tau}
\end{table}

\subsection{The net rate: what if another team uses a safety net too?}
\label{sec:net}
Left alone, the optimizer picks the highest allowed net (0.989), because none of the twelve scenarios contains a
competitor net below 0.99. But a net is only taken by applicants nobody else serves more cheaply, so its level
matters only \emph{relative to other nets}, and the 2\% preference makes that comparison a step function: against
a competitor net at $x$ we win every abandoned applicant if ours is below $x - 0.02$, only our own fans if it is
within $\pm 0.02$ of $x$, and nothing if it is above $x + 0.02$.

\paragraph{Step 1: competitors who net above a high threshold.} We first let the base competitors (who reject
above a 25\% PD) replace ``no offer'' by a net at 1.00, 0.99, 0.95, 0.90, 0.85, 0.80 or 0.70, thirteen cases in
all including non-netting ones. Lowering our net from 0.989 to 0.90 costs 0.14~\meur{} if nobody else nets and
gains 1.4~\meur{} on average if somebody does; below 0.85 the lost margin outweighs the protection. Cheap
insurance, but small stakes: these competitors abandon only 0.8\% of the applicants.

\paragraph{Step 2: clones.} A team that reasons like us would not only net; it would also stop lending at
normal rates in the 10--25\% PD band, as we do. We call such a team a \emph{clone}: it applies our rule
(Table~\ref{tab:rule}) to its own PD estimate, with a net of its choosing. The consequences are of a different
order of magnitude:
\begin{itemize}
  \item With two clones, the band where all three lenders net is \textbf{34\% of applicants} (73\% of the ones we
        net). It defaults at 18.7\%, so its break-even net rate is only 0.23 and a loan at 83\% earns about
        5,000~\texteuro{} in expectation.
  \item Whoever is lowest within 2\% takes the band: 34,000 loans, \textbf{65~\meur{}} if we only keep our own
        fans, \textbf{150 to 190~\meur{}} if we take all of it. The whole normal game is worth 5~\meur{}.
  \item Hence, if any team nets the band, the ranking of stage 1 is decided by the net war, not by careful
        pricing. The rules make this possible because borrowers have no outside option: whoever receives an offer
        borrows, at 83\% if need be.
\end{itemize}

\begin{table}[h]\centering\footnotesize
\begin{tabular}{lrr@{\hskip 0.8cm}rrrrrr}
\toprule
 & Nobody & One clone at & \multicolumn{6}{c}{Both competitors are clones with a net at}\\
Our net & nets & 0.898, one base & 0.989 & 0.95 & 0.898 & 0.87 & 0.85 & 0.80\\
\midrule
0.989 & 6.1 &  4.8 &  75 &   5 &   5 &   5 &   5 &   5\\
0.898 & 5.7 & 13.3 & 190 & 190 &  67 &   5 &   5 &   5\\
0.877 & 5.6 & 16.9 & 184 & 184 & 184 &  65 &   5 &   5\\
0.830 & 5.4 & 15.6 & 171 & 171 & 171 & 171 & 171 &   5\\
0.750 & 5.0 & 13.4 & 149 & 149 & 149 & 149 & 149 & 149\\
\bottomrule
\end{tabular}
\caption{Our expected profit (\meur{}) by the level of our net, against clones at various nets (margins and
thresholds as chosen). ``Nobody nets'' is the base scenario.}
\label{tab:clones}
\end{table}

Table~\ref{tab:clones} shows the step structure at work. At 0.898 we would take everything from a clone at 0.95
or above, keep only our fans against a clone at exactly 0.898 (an exact tie is plausible: it is the obvious
``just below 0.90''), and get \emph{nothing} from a clone at 0.87 or 0.85 -- a regret of 100 to 170~\meur{}.
Going lower is almost free in the world where nobody nets: from 0.898 to 0.83 costs 0.33~\meur{}.

\paragraph{The decision.} Table~\ref{tab:decision} weighs 26 competitor cases (nobody nets; base competitors
with nets; moderate clones that keep lending normally up to twice our thresholds; one clone and one base
competitor; two clones) under three priors on how likely clones are. Within a family the cases are
equally likely; the family probabilities are: skeptical 70 / 15 / 5 / 7 / 3\%, moderate 50 / 20 / 10 / 12 / 8\%,
paranoid 30 / 20 / 15 / 15 / 20\%.

\begin{table}[h]\centering\small
\begin{tabular}{lrrr}
\toprule
Our net & Skeptical prior & Moderate prior & Paranoid prior\\
\midrule
0.989 &  5.4 &  5.5 &  6.1\\
0.950 &  6.0 &  7.1 &  9.6\\
0.920 &  6.6 &  8.8 & 13.6\\
0.898 &  7.3 & 10.5 & 17.5\\
0.877 &  8.0 & 12.3 & 22.0\\
0.850 &  8.6 & 14.0 & 25.9\\
\textbf{0.830} & \textbf{8.8} & \textbf{14.5} & \textbf{27.2}\\
0.800 &  8.6 & 14.2 & 26.8\\
0.750 &  8.6 & 14.3 & 27.1\\
0.700 &  8.4 & 14.1 & 26.8\\
0.600 &  7.4 & 12.2 & 23.3\\
No offer & 4.7 & 4.7 & 4.8\\
\bottomrule
\end{tabular}
\caption{Prior-weighted expected profit (\meur{}) as a function of our net rate, margins and thresholds as
chosen.}
\label{tab:decision}
\end{table}

Three conclusions follow.
\begin{enumerate}
  \item \textbf{0.83 is the best level under every prior}, even the skeptical one where both competitors are
        clones with only 3\% probability, and the curve is flat between 0.75 and 0.85. The asymmetry drives it:
        being too high against a clone costs 100~\meur{}, being too low when nobody nets costs 0.3~\meur{}.
        Below 0.70 the net starts winning genuinely terrible borrowers from competitors who lend to everyone at
        70--85\%, which sets the floor.
  \item \textbf{Lowering the thresholds further does not pay.} Re-optimizing the whole rule with clones in the
        scenario mix nets even more applicants (above a 4--8\% PD at low income) with a net at 0.70. That rule
        earns 1.1~\meur{} less in the world where nobody nets (4.1 against 5.2~\meur{}) and only beats the chosen
        rule under the paranoid prior. We keep the thresholds of Table~\ref{tab:rule} and will revisit them in
        stage 2 with real information.
  \item \textbf{One tick below the round number.} At exactly $x - 0.02$ we tie on the clone's own fans, so we
        submit \textbf{0.828} rather than 0.83: it beats a competitor net at 0.85, 0.898, 0.90, 0.95, 0.99 or 1.00
        for all three applicant types, splits with a net between 0.81 and 0.848, and gets nothing from a net
        below 0.808. Its cost if nobody nets at all is 0.35~\meur{}, about 6\% of our profit in that world.
\end{enumerate}

\section{The submission}
\code{output/round1/rates\_round1.csv}: 100,000 rows, columns \code{id} (1 to 100,000, in order) and \code{rate},
no missing value, maximum 0.828. 54,418 normal offers (7,270 below 4\%; 13,250 at 4--6\%; 11,382 at 6--8\%;
14,389 at 8--12\%; 8,127 at 12--17.6\%) and 45,582 offers at 0.828. No rate is a round number, to avoid ties with
competitors who use them.

\section{Assumptions, risks and what could go wrong}
\begin{itemize}
  \item \textbf{Competitor behaviour is a guess.} The archetypes span cautious to aggressive and include the
        lecture-formula mistake, selection-aware teams and clones of ourselves, and we optimized the average while
        checking the worst case; but a behaviour outside this family (for instance random rates, or a flat 10\%
        for everyone) was not simulated. The rule stays profitable in every scenario we tried; the main risk is
        \emph{lower} profit, not a loss.
  \item \textbf{The net war is winner-take-all.} A clone netting below 0.808 takes the whole band and leaves us
        5~\meur{} against its 150~\meur{} or more. No choice of ours is safe against every level; 0.828 is a bet
        that a competitor's net, if any, sits at 0.85 or above (the natural choices being 1.00, 0.99, 0.95, 0.90
        and 0.85).
  \item \textbf{The net earns nothing if the competitors lend to everybody}, and nothing if they both net below
        0.808; neither case costs us money. It earns 0.4~\meur{} in the base scenario and far more if they reject
        at 20\% or below. Against non-clone rejecters its value rests on the abandoned pool defaulting below 45\%,
        which holds for any rejection threshold up to 25\% and fails only for a 0.2\% pool.
  \item \textbf{The rules give borrowers no outside option.} If the simulation applies an unstated reservation
        rate, the net earns nothing and costs nothing. The tie-breaking rule is also unknown; we avoided round
        numbers.
  \item \textbf{Our thresholds give up share.} We lend to 9--22\% of applicants depending on the scenario; above
        10k of income we win our full third, below we deliberately walk away from the risky middle. This is a
        choice, not an accident.
  \item \textbf{The competitors' models are assumed as good as ours.} If theirs are worse, the curse is smaller
        and their rejections noisier (better abandoned pool); if better, the curse is larger and the low thresholds
        protect us. Applicant types are assumed uniform and independent of characteristics, as stated in the rules.
  \item \textbf{Visibility.} 45\% of our offers at 82.8\% will be noticed in the stage-1 results. If it works, the
        other teams will copy it in stage~2 and the band will be competed down towards its break-even. We bank it
        while it is unique.
\end{itemize}

\section{Plan for stage 2}
The stage-1 results give us the realized market: who each applicant borrowed from, who defaulted, and -- we hope
-- the competitors' rates. The first thing to read is the net war: which teams netted, at what level and with
what thresholds, and who took the band; that alone will explain most of the profit ranking. Then (i) infer each
competitor's actual pricing function (regressing their rate on $X$ and the signals if rates are disclosed,
otherwise from who won which applicant), (ii) measure the realized curse on our own loans segment by segment and
compare it with the simulation, (iii) replace the archetypes with the two estimated competitors and re-optimize,
and (iv) only then consider aggressive moves: undercutting where a competitor is shown to overprice, moving the
net against theirs, revisiting the thresholds, and re-opening the 10--25\% PD band if the measured curse turns
out smaller than simulated. Everything in this note is reproducible, so stage 2 is a change of inputs, not a
rewrite.

\section{Files and how to reproduce}
Everything lives in the project folder.
\begingroup\sloppy
\begin{itemize}
  \item \code{run\_round1.py} -- the whole pipeline. It regenerates the submission in about two minutes and
        prints every table of Sections 4--5.5:
        \begin{center}\code{python run\_round1.py --no-sex --net 0.828}\end{center}
  \item \code{lending\_game/} -- the package: \code{features.py}, \code{models.py}, \code{simulator.py}
        (rules of the game), \code{competitors.py} (archetypes, scenarios, clones), \code{strategy.py} (our rule,
        optimizer).
  \item \code{analysis/} -- \code{robustness\_round1.py} (thresholds, scenario weighting),
        \code{net\_rate\_scenarios.py} (Step 1 of Section~\ref{sec:net}) and \code{clone\_scenarios.py}
        (Step 2 and the decision table).
  \item \code{output/round1/}:
    \begin{itemize}
      \item \code{rates\_round1.csv} -- the submission;
      \item \code{strategy\_round1.json} -- all parameters of the rule and of the model;
      \item \code{diagnostics\_round1.md} -- every table printed by the pipeline;
      \item \code{clone\_profit\_table\_*.csv} -- our profit by net level against every clone case;
      \item \code{new\_applications\_scored.csv} -- our PD and rate per applicant, to diagnose stage 1 later;
      \item \code{oos\_probas\_sex0.npz} -- out-of-sample PDs of the four models on the past loans.
    \end{itemize}
\end{itemize}
\endgroup

\end{document}
